\chapter{2007年上海卷（春）}

\section{填空题}

\begin{problem}
计算 \( \displaystyle\lim_{n\to\infty}\frac{2n^{2}+1}{3n(n+1)}= \)\fillinblank{}.
\end{problem}

\begin{answer}
\(\frac{2}{3}\)
\end{answer}

\begin{problem}
若关于 \(x\) 的一元二次实系数方程 \(x^{2}+px+q=0\) 有一个根为 \(1+\i\)（\(\i\) 是虚数单位），则 \(q=\)\fillinblank{}.
\end{problem}

\begin{answer}
\(2\)
\end{answer}

\begin{problem}
若关于 \(x\) 的不等式 \(\frac{x-a}{x+1}>0\) 的解集为 \((-\infty,-1)\cup(4,+\infty)\)，则实数 \(a=\)\fillinblank{}.
\end{problem}

\begin{answer}
\(4\)
\end{answer}

\begin{problem}
函数 \(y=(\sin x+\cos x)^2\) 的最小正周期为\fillinblank{}.
\end{problem}

\begin{answer}
\(\pi\)
\end{answer}

\begin{problem}
设函数 \(y=f(x)\) 是奇函数. 若 \(f(-2)+f(-1)-3=f(1)+f(2)+3\)，则 \(f(1)+f(2)=\)\fillinblank{}.
\end{problem}

\begin{answer}
\(-3\)
\end{answer}

\begin{problem}
在平面直角坐标系 \(xOy\) 中，若抛物线 \(y^2=4x\) 上的点 \(P\) 到该抛物线的焦点的距离为 \(6\)，则点 \(P\) 的横坐标 \(x=\)\fillinblank{}.
\end{problem}

\begin{answer}
\(5\)
\end{answer}

\begin{problem}
在平面直角坐标系 \(xOy\) 中，若曲线 \(x=\sqrt{4-y^2}\) 与直线 \(x=m\) 有且只有一个公共点，则实数 \(m=\)\fillinblank{}.
\end{problem}

\begin{answer}
\(2\)
\end{answer}

\begin{problem}
若向量 \(\bs{a}\)，\(\bs{b}\) 满足 \(|\bs{a}|=\sqrt{2}\)，\(|\bs{b}|=1\)，\(\bs{a}\cdot(\bs{a}+\bs{b})=1\)，则向量 \(\bs{a}\)，\(\bs{b}\) 的夹角的大小为\fillinblank{}.
\end{problem}

\begin{answer}
\(\frac{3\pi}{4}\)
\end{answer}

\begin{problem}
若 \(x_1\)，\(x_2\) 为方程 \(2^x=\left(\frac{1}{2}\right)^{-\frac{1}{x}+1}\) 的两个实数解，则 \(x_1+x_2=\)\fillinblank{}.
\end{problem}

\begin{answer}
\(-1\)
\end{answer}

\begin{problem}
在一次教师联欢会上，到会的女教师比男教师多 \(12\) 人，从这些教师中随机挑选一人表演节目. 若选到男教师的概率为 \(\frac{9}{20}\)，则参加联欢会的教师共有\fillinblank{}人.
\end{problem}

\begin{answer}
\(120\)
\end{answer}

\begin{problem}
函数 \(y=
\begin{cases}
x^2+1, & x\geqslant 0,\\
\frac{2}{x}, & x<0
\end{cases}
\)
的反函数是\fillinblank{}.
\end{problem}

\begin{answer}
\(y=
\begin{cases}
\sqrt{x-1}, & x\geqslant 1,\\
\frac{2}{x}, & x<0
\end{cases}
\)
\end{answer}

\section{单选题}

\begin{problem}
若集合 \(A=\{1,m^2\}\)，\(B=\{2,4\}\)，则“\(m=2\)”是“\(A\cap B=\{4\}\)”的（\quad）

\choices
    {充分不必要条件.}
    {必要不充分条件.}
    {充要条件.}
    {既不充分也不必要条件.}
\end{problem}

\begin{answer}
A
\end{answer}

\begin{problem}
如图，平面内的两条相交直线 \(OP_1\) 和 \(OP_2\) 将该平面分割成四个部分 \(\mathrm{I}\)，\(\mathrm{II}\)，\(\mathrm{III}\)，\(\mathrm{IV}\)（不包括边界）. 若 \(\overrightarrow{OP}=a\overrightarrow{OP_1}+b\overrightarrow{OP_2}\)，且点 \(P\) 落在第 \(\mathrm{III}\) 部分，则实数 \(a\)，\(b\) 满足（\quad）

\bitmapfigure[width=0.30\linewidth]{img/2007/shanghai_spring/q13_fig1.png}

\choices
    {\(a>0\)，\(b>0\).}
    {\(a>0\)，\(b<0\).}
    {\(a<0\)，\(b>0\).}
    {\(a<0\)，\(b<0\).}
\end{problem}

\begin{answer}
B
\end{answer}

\begin{problem}
下列四个函数中，图象如图所示的只能是（\quad）

\bitmapfigure[width=0.24\linewidth]{img/2007/shanghai_spring/q14_fig1.png}

\choices
    {\(y=x+\lg x\).}
    {\(y=x-\lg x\).}
    {\(y=-x+\lg x\).}
    {\(y=-x-\lg x\).}
\end{problem}

\begin{answer}
B
\end{answer}

\begin{problem}
设 \(a\)，\(b\) 是正实数，以下不等式：

\begin{circlelist}
\item \(\sqrt{ab}>\frac{2ab}{a+b}\)；
\item \(a>|a-b|-b\)；
\item \(a^2+b^2>4ab-3b^2\)；
\item \(ab+\frac{2}{ab}>2\).
\end{circlelist}

恒成立的序号为（\quad）

\choices
    {\(\circled{1}\)，\(\circled{3}\).}
    {\(\circled{1}\)，\(\circled{4}\).}
    {\(\circled{2}\)，\(\circled{3}\).}
    {\(\circled{2}\)，\(\circled{4}\).}
\end{problem}

\begin{answer}
D
\end{answer}

\section{解答题}

\begin{problem}
如图，在棱长为 \(2\) 的正方体 \(ABCD-A'B'C'D'\) 中，点 \(E\)，\(F\) 分别是 \(A'B'\) 和 \(AB\) 的中点，求异面直线 \(A'F\) 与 \(CE\) 所成角的大小（结果用反三角函数值表示）.

\bitmapfigure[width=0.36\linewidth]{img/2007/shanghai_spring/q16_fig1.png}
\end{problem}

\begin{solution}
【法一】如图建立空间直角坐标系. 由题意可知
\(A'(2,0,2)\)，\(C(0,2,0)\)，\(E(2,1,2)\)，\(F(2,1,0)\)，\(\overrightarrow{A'F}=(0,1,-2)\)，\(\overrightarrow{CE}=(2,-1,2)\).
设直线 \(A'F\) 与 \(CE\) 所成角为 \(\theta\)，则
\examdisplaymath{
\cos\theta=
\frac{\left|\overrightarrow{A'F}\cdot\overrightarrow{CE}\right|}{\left|\overrightarrow{A'F}\right|\cdot\left|\overrightarrow{CE}\right|}
=\frac{5}{\sqrt{5}\cdot 3}
=\frac{\sqrt{5}}{3}
}
因此 \(\theta=\arccos\frac{\sqrt{5}}{3}\)，即异面直线 \(A'F\) 与 \(CE\) 所成角的大小为 \(\arccos\frac{\sqrt{5}}{3}\).

【法二】连结 \(EB\). 因为 \(A'E\parallel BF\)，且 \(A'E=BF\)，所以 \(A'FBE\) 是平行四边形，从而 \(A'F\parallel EB\). 因此异面直线 \(A'F\) 与 \(CE\) 所成的角就是 \(CE\) 与 \(EB\) 所成的角. 又因为 \(CB\perp\) 平面 \(ABB'A'\)，所以 \(CB\perp BE\). 在直角三角形 \(CEB\) 中，\(CB=2\)，\(BE=\sqrt{5}\)，于是
\(\tan\angle CEB=\frac{2\sqrt{5}}{5}\).
故 \(\angle CEB=\arctan\frac{2\sqrt{5}}{5}\). 因此异面直线 \(A'F\) 与 \(CE\) 所成角的大小为 \(\arctan\frac{2\sqrt{5}}{5}\).
\end{solution}

\begin{problem}
求出一个数学问题的正确结论后，将其作为条件之一，提出与原来问题有关的新问题，我们把它称为原来问题的一个“逆向”问题. 例如，原来问题是“若正四棱锥底面边长为 \(4\)，侧棱长为 \(3\)，求该正四棱锥的体积”. 求出体积 \(\frac{16}{3}\) 后，它的一个“逆向”问题可以是“若正四棱锥底面边长为 \(4\)，体积为 \(\frac{16}{3}\)，求侧棱长”；也可以是“若正四棱锥的体积为 \(\frac{16}{3}\)，求所有侧面面积之和的最小值”. 试给出问题“在平面直角坐标系 \(xOy\) 中，求点 \(P(2,1)\) 到直线 \(3x+4y=0\) 的距离.”的一个有意义的“逆向”问题，并解答你所给出的“逆向”问题.
\end{problem}

\begin{solution}
点 \(P(2,1)\) 到直线 \(3x+4y=0\) 的距离为 \(\frac{|3\cdot 2+4\cdot 1|}{\sqrt{3^2+4^2}}=2\).

“逆向”问题可以是：

求到直线 \(3x+4y=0\) 的距离为 \(2\) 的点的轨迹方程. 设所求轨迹上任意一点为 \(P(x,y)\)，则 \(\frac{|3x+4y|}{5}=2\)，所求轨迹为 \(3x+4y-10=0\) 或 \(3x+4y+10=0\).

若点 \(P(2,1)\) 到直线 \(l:ax+by=0\) 的距离为 \(2\)，求直线 \(l\) 的方程. 由 \(\frac{|2a+b|}{\sqrt{a^2+b^2}}=2\)，化简得 \(4ab-3b^2=0\)，即 \(b=0\) 或 \(4a=3b\). 所以直线 \(l\) 的方程为 \(x=0\) 或 \(3x+4y=0\).

意义不大的“逆向”问题可能是：

点 \(P(2,1)\) 是不是到直线 \(3x+4y=0\) 的距离为 \(2\) 的一个点？因为 \(\frac{|3\cdot2+4\cdot1|}{\sqrt{3^2+4^2}}=2\)，所以点 \(P(2,1)\) 是到直线 \(3x+4y=0\) 的距离为 \(2\) 的一个点.

点 \(Q(1,1)\) 是不是到直线 \(3x+4y=0\) 的距离为 \(2\) 的一个点？因为 \(\frac{|3\cdot1+4\cdot1|}{\sqrt{3^2+4^2}}=\frac{7}{5}\neq 2\)，所以点 \(Q(1,1)\) 不是到直线 \(3x+4y=0\) 的距离为 \(2\) 的一个点.

点 \(P(2,1)\) 是不是到直线 \(5x+12y=0\) 的距离为 \(2\) 的一个点？因为 \(\frac{|5\cdot2+12\cdot1|}{\sqrt{5^2+12^2}}=\frac{22}{13}\neq 2\)，所以点 \(P(2,1)\) 不是到直线 \(5x+12y=0\) 的距离为 \(2\) 的一个点.
\end{solution}

\begin{problem}
如图，在直角坐标系 \(xOy\) 中，设椭圆 \(C:\frac{x^2}{a^2}+\frac{y^2}{b^2}=1\)（\(a>b>0\)）的左右两个焦点分别为 \(F_1\)，\(F_2\). 过右焦点 \(F_2\) 且与 \(x\) 轴垂直的直线 \(l\) 与椭圆 \(C\) 相交，其中一个交点为 \(M(\sqrt{2},1)\).

\begin{enumerate}
\item 求椭圆 \(C\) 的方程；
\item 设椭圆 \(C\) 的一个顶点为 \(B(0,-b)\)，直线 \(BF_2\) 交椭圆 \(C\) 于另一点 \(N\)，求 \(\triangle F_1BN\) 的面积.
\end{enumerate}

\bitmapfigure[width=0.38\linewidth]{img/2007/shanghai_spring/q18_fig1.png}
\end{problem}

\begin{solution}
\begin{enumerate}
\item 【法一】因为 \(l\) 垂直于 \(x\) 轴，所以 \(F_2\) 的坐标为 \((\sqrt{2},0)\). 由题意可知 \(\frac{2}{a^2}+\frac{1}{b^2}=1\)，\(a^2-b^2=2\).
解得 \(a^2=4\)，\(b^2=2\). 因此所求椭圆方程为 \(\frac{x^2}{4}+\frac{y^2}{2}=1\).

【法二】由椭圆定义可知 \(|MF_1|+|MF_2|=2a\). 由题意 \(|MF_2|=1\)，因此 \(|MF_1|=2a-1\). 又由直角三角形 \(MF_1F_2\) 可知 \((2a-1)^2=(2\sqrt{2})^2+1\)，且 \(a>0\)，所以 \(a=2\). 再由 \(a^2-b^2=2\) 得 \(b^2=2\). 因此椭圆 \(C\) 的方程为 \(\frac{x^2}{4}+\frac{y^2}{2}=1\).

\item 直线 \(BF_2\) 的方程为 \(y=x-\sqrt{2}\). 由 \(y=x-\sqrt{2}\)，\(\frac{x^2}{4}+\frac{y^2}{2}=1\)
得点 \(N\) 的纵坐标为 \(\frac{\sqrt{2}}{3}\). 又 \(|F_1F_2|=2\sqrt{2}\)，所以 \(S_{\triangle F_1BN}=\frac{1}{2}\times\left(\sqrt{2}+\frac{\sqrt{2}}{3}\right)\times 2\sqrt{2}=\frac{8}{3}\).
\end{enumerate}
\end{solution}

\begin{problem}
某人定制了一批地砖. 每块地砖（如图 1 所示）是边长为 \(0.4\text{ 米}\) 的正方形 \(ABCD\)，点 \(E\)，\(F\) 分别在边 \(BC\) 和 \(CD\) 上，\(\triangle CFE\)，\(\triangle ABE\) 和四边形 \(AEFD\) 均由单一材料制成，制成 \(\triangle CFE\)，\(\triangle ABE\) 和四边形 \(AEFD\) 的三种材料的每平方米价格之比依次为 \(3:2:1\). 若将此种地砖按图 2 所示的形式铺设，能使中间的深色阴影部分成四边形 \(EFGH\).

\begin{enumerate}
\item 求证：四边形 \(EFGH\) 是正方形；
\item \(E\)，\(F\) 在什么位置时，定制这批地砖所需的材料费用最省？
\end{enumerate}

\examfiguregroup[float=inline]{img/2007/shanghai_spring/q19_fig1.png}{%
    \bitmapinclude[width=0.28\linewidth]{img/2007/shanghai_spring/q19_fig1.png}\hfill
    \bitmapinclude[width=0.46\linewidth]{img/2007/shanghai_spring/q19_fig2.png}%
}
\end{problem}

\begin{solution}
\begin{enumerate}
\item 图 2 是由四块图 1 所示地砖绕点 \(C\) 按顺时针旋转 \(90^\circ\) 后得到，\(\triangle CFE\) 为等腰直角三角形，所以四边形 \(EFGH\) 是正方形.

\item 设 \(CE=x\)，则 \(BE=0.4-x\). 每块地砖的费用为 \(W\). 制成 \(\triangle CFE\)，\(\triangle ABE\) 和四边形 \(AEFD\) 三种材料的每平方米价格依次为 \(3a\)，\(2a\)，\(a\)（元），于是
\examdisplaymath{
\begin{aligned}
W&=\frac{1}{2}x^2\cdot 3a+\frac{1}{2}\times 0.4\times(0.4-x)\times 2a\\
&\quad+\left[0.16-\frac{1}{2}x^2-\frac{1}{2}\times 0.4\times(0.4-x)\right]a\\
&=a(x^2-0.2x+0.24)
=a\left[(x-0.1)^2+0.23\right],\quad 0<x<0.4.
\end{aligned}
}
由 \(a>0\) 可知，当 \(x=0.1\) 时，\(W\) 有最小值，即总费用最省. 答：当 \(CE=CF=0.1\text{ 米}\) 时，总费用最省.
\end{enumerate}
\end{solution}

\begin{problem}
通常用 \(a\)，\(b\)，\(c\) 分别表示 \(\triangle ABC\) 的三个内角 \(A\)，\(B\)，\(C\) 所对边的边长，\(R\) 表示 \(\triangle ABC\) 的外接圆半径.

\begin{enumerate}
\item 如图，在以 \(O\) 为圆心，半径为 \(2\) 的 \(\odot O\) 中，\(BC\) 和 \(BA\) 是 \(\odot O\) 的弦，其中 \(BC=2\)，\(\angle ABC=45^\circ\)，求弦 \(AB\) 的长；
\item 在 \(\triangle ABC\) 中，若 \(\angle C\) 是钝角，求证：\(a^2+b^2<4R^2\)；
\item 给定三个正实数 \(a\)，\(b\)，\(R\)，其中 \(b\leqslant a\). 问：\(a\)，\(b\)，\(R\) 满足怎样的关系时，以 \(a\)，\(b\) 为边长，\(R\) 为外接圆半径的 \(\triangle ABC\) 不存在，存在一个或存在两个（全等三角形算作同一个）？在 \(\triangle ABC\) 存在的情况下，用 \(a\)，\(b\)，\(R\) 表示 \(c\).
\end{enumerate}

\bitmapfigure[width=0.28\linewidth]{img/2007/shanghai_spring/q20_fig1.png}
\end{problem}

\begin{solution}
\begin{enumerate}
\item \(\triangle ABC\) 的外接圆半径为 \(2\). 在 \(\triangle ABC\) 中，\(AC=2R\sin B=2\sqrt{2}\)，\(\sin A=\frac{BC}{2R}=\frac{1}{2}\)，所以 \(A=30^\circ\). 因而
\examdisplaymath{
AB^2=BC^2+AC^2-2BC\cdot AC\cdot\cos C
=4+8+8\sqrt{2}\cos(A+B)=4(\sqrt{3}+2)=2(\sqrt{3}+1)^2
}
所以 \(AB=\sqrt{6}+\sqrt{2}\).

\item 由 \(\sin A=\frac{a}{2R}\)，\(\sin B=\frac{b}{2R}\). 因为 \(\angle C\) 是钝角，所以 \(\angle A\)，\(\angle B\) 都是锐角，于是
\(\cos A=\frac{1}{2R}\sqrt{4R^2-a^2}\)，\(\cos B=\frac{1}{2R}\sqrt{4R^2-b^2}\)，且
\examdisplaymath{
\cos C=-\cos(A+B)=\frac{1}{4R^2}\left(ab-\sqrt{4R^2-a^2}\sqrt{4R^2-b^2}\right)<0
}
因此 \(a^2b^2<(4R^2-a^2)(4R^2-b^2)\)，从而 \(16R^4-4R^2(a^2+b^2)>0\). 即 \(a^2+b^2<4R^2\).

\item 当 \(a>2R\) 或 \(a=b=2R\) 时，所求的 \(\triangle ABC\) 不存在.

当 \(a=2R\) 且 \(b<a\) 时，\(\angle A=90^\circ\)，所求的 \(\triangle ABC\) 只存在一个，且 \(c=\sqrt{a^2-b^2}\).

当 \(a<2R\) 且 \(b=a\) 时，\(\angle A=\angle B\)，且 \(A\)，\(B\) 都是锐角. 由 \(\sin A=\frac{a}{2R}=\frac{b}{2R}=\sin B\)，可知 \(A\)，\(B\) 唯一确定. 因此所求的 \(\triangle ABC\) 只存在一个，且
\(c=2a\cos A=\frac{a}{R}\sqrt{4R^2-a^2}\).

当 \(b<a<2R\) 时，\(\angle B\) 总是锐角，\(\angle A\) 可以是钝角，也可以是锐角，因此所求的 \(\triangle ABC\) 存在两个. 由 \(\sin A=\frac{a}{2R}\)，\(\sin B=\frac{b}{2R}\) 得：

若 \(\angle A<90^\circ\)，则 \(\cos A=\frac{1}{2R}\sqrt{4R^2-a^2}\)，
\examdisplaymath{
c=\sqrt{a^2+b^2+\frac{ab}{2R^2}\left(\sqrt{4R^2-a^2}\sqrt{4R^2-b^2}-ab\right)}
}

若 \(\angle A>90^\circ\)，则 \(\cos A=-\frac{1}{2R}\sqrt{4R^2-a^2}\)，
\examdisplaymath{
c=\sqrt{a^2+b^2-\frac{ab}{2R^2}\left(\sqrt{4R^2-a^2}\sqrt{4R^2-b^2}+ab\right)}
}
\end{enumerate}
\end{solution}

\begin{problem}
我们在下面的表格内填写数值：先将第 1 行的所有空格填上 \(1\)；再把一个首项为 \(1\)，公比为 \(q\) 的数列 \(\{a_n\}\) 依次填入第一列的空格内；然后按照“任意一格的数是它上面一格的数与它左边一格的数之和”的规则填写其他空格.

\begin{center}
\begin{tabular}{|c|c|c|c|c|c|}
\hline
 & 第 1 列 & 第 2 列 & 第 3 列 & \(\cdots\) & 第 \(n\) 列 \\
\hline
第 1 行 & \(1\) & \(1\) & \(1\) & \(\cdots\) & \(1\) \\
\hline
第 2 行 & \(q\) &  &  &  &  \\
\hline
第 3 行 & \(q^2\) &  &  &  &  \\
\hline
\(\cdots\) & \(\cdots\) &  &  &  &  \\
\hline
第 \(n\) 行 & \(q^{n-1}\) &  &  &  &  \\
\hline
\end{tabular}
\end{center}

\begin{enumerate}
\item 设第 2 行的数依次为 \(B_1\)，\(B_2\)，\(\cdots\)，\(B_n\)，试用 \(n\)，\(q\) 表示 \(B_1+B_2+\cdots+B_n\) 的值；
\item 设第 3 列的数依次为 \(c_1\)，\(c_2\)，\(c_3\)，\(\cdots\)，\(c_n\)，求证：对于任意非零实数 \(q\)，\(c_1+c_3>2c_2\)；
\item 请在以下两个问题中选择一个进行研究（只能选择一个问题，如果都选，被认为选择了第一问）：
\begin{enumerate}
\item 能否找到 \(q\) 的值，使得第二问中的数列 \(c_1\)，\(c_2\)，\(c_3\)，\(\cdots\)，\(c_n\) 的前 \(m\) 项 \(c_1\)，\(c_2\)，\(c_3\)，\(\cdots\)，\(c_m\)（\(m\geqslant3\)）成为等比数列？若能找到，\(m\) 的值有多少个？若不能找到，说明理由；
\item 能否找到 \(q\) 的值，使得填完表格后，除第 1 列外，还有不同的两列数的前三项各自依次成等比数列？并说明理由.
\end{enumerate}
\end{enumerate}
\end{problem}

\begin{solution}
\begin{enumerate}
\item \(B_1=q\)，\(B_2=1+q\)，\(B_3=1+(1+q)=2+q\)，\(\cdots\)，\(B_n=(n-1)+q\). 所以 \(B_1+B_2+\cdots+B_n=1+2+\cdots+(n-1)+nq=\frac{n(n-1)}{2}+nq\).

\item \(c_1=1\)，\(c_2=1+(1+q)=2+q\)，\(c_3=(2+q)+(1+q+q^2)=3+2q+q^2\). 因而
\examdisplaymath{
c_1+c_3-2c_2=1+3+2q+q^2-2(2+q)=q^2>0
}
所以 \(c_1+c_3>2c_2\).

\item
\begin{enumerate}
\item 先设 \(c_1\)，\(c_2\)，\(c_3\) 成等比数列. 由 \(c_1c_3=c_2^2\)，得 \(3+2q+q^2=(2+q)^2\)，所以 \(q=-\frac{1}{2}\). 此时 \(c_1=1\)，\(c_2=\frac{3}{2}\)，\(c_3=\frac{9}{4}\)，所以 \(c_1\)，\(c_2\)，\(c_3\) 是一个公比为 \(\frac{3}{2}\) 的等比数列. 如果 \(m\geqslant4\) 时 \(c_1\)，\(c_2\)，\(\cdots\)，\(c_m\) 为等比数列，那么 \(c_1\)，\(c_2\)，\(c_3\) 一定是等比数列. 由上所述，此时 \(q=-\frac{1}{2}\)，且 \(c_4=\frac{23}{8}\). 因为 \(\frac{c_4}{c_3}\neq\frac{3}{2}\)，所以对于任意 \(m\geqslant4\)，数列 \(c_1\)，\(c_2\)，\(\cdots\)，\(c_m\) 一定不是等比数列. 综上所述，当且仅当 \(m=3\) 且 \(q=-\frac{1}{2}\) 时，数列 \(c_1\)，\(c_2\)，\(\cdots\)，\(c_m\) 是等比数列.

\item 设 \(x_1\)，\(x_2\)，\(x_3\) 和 \(y_1\)，\(y_2\)，\(y_3\) 分别是第 \(k+1\) 列和第 \(m+1\) 列的前三项，其中 \(1\leqslant k<m\leqslant n-1\). 则
\(x_1=1\)，\(x_2=k+q\)，\(x_3=\frac{k(k+1)}{2}+kq+q^2\).
若第 \(k+1\) 列的前三项 \(x_1\)，\(x_2\)，\(x_3\) 是等比数列，则由 \(x_1x_3=x_2^2\) 得 \(\frac{k(k+1)}{2}+kq+q^2=(k+q)^2\)，
从而 \(\frac{k^2-k}{2}+kq=0\)，所以 \(q=\frac{1-k}{2}\). 同理，若第 \(m+1\) 列的前三项 \(y_1\)，\(y_2\)，\(y_3\) 是等比数列，则 \(q=\frac{1-m}{2}\). 当 \(k\neq m\) 时，\(\frac{1-k}{2}\neq\frac{1-m}{2}\). 所以无论怎样的 \(q\)，都不能同时找到两列数（除第 1 列外），使它们的前三项都成等比数列.
\end{enumerate}
\end{enumerate}
\end{solution}
